\documentclass[12pt, a4paper]{article}

% --- Encoding & Fonts ---
\usepackage[utf8]{inputenc}
\usepackage[T1]{fontenc}
\usepackage{lmodern}

% --- Mathematics ---
\usepackage{amsmath}
\usepackage{amssymb}
\usepackage{amsthm}
\usepackage{mathtools}

% --- Layout ---
\usepackage[margin=1in]{geometry}

% --- Lists ---
\usepackage{enumitem}

% --- Tables ---
\usepackage{booktabs}
\usepackage{array}

% --- Graphics ---
\usepackage{graphicx}

% --- Headers / Footers ---
\usepackage{fancyhdr}
\pagestyle{fancy}
\fancyhf{}
\fancyhead[L]{\leftmark}
\fancyhead[R]{\thepage}
\renewcommand{\headrulewidth}{0.4pt}

% --- Hyperlinks & Cross-references ---
\usepackage[colorlinks=true, linkcolor=blue, urlcolor=blue, citecolor=blue]{hyperref}

% --- Misc ---
\usepackage{xcolor}

% --- Custom commands (safety definitions) ---
\providecommand{\operatorname}[1]{\mathop{\mathrm{#1}}\nolimits}

\title{Learning Module: Sdos2}
\author{}
\date{\today}

\begin{document}

\maketitle
\thispagestyle{empty}
\newpage

\tableofcontents
\newpage

% ============================================================
% SECTION  Working in Teams
% ============================================================

% !!!  REVIEW FLAG: Working in Teams  !!!
% [knowledge_assumption] Add a brief explanatory note or minilesson on calculating percentagepoint increases in earlier sections or embed a brief calculation example in the question stem.
% !!! END REVIEW FLAGS !!!
\section{Working in Teams}

\subsection{Motivation}
Imagine a soccer team where each player focuses only on scoring versus working together to win the match; the latter reflects a highperforming software team.  
Software projects succeed not because a single ``genius programmer'' writes perfect code, but because the whole team works effectively toward shared goals.  
Understanding what makes a team effective helps prevent project failures that arise from poor collaboration, low motivation, or dissatisfied users.

\subsection{Intuition}
Think of a software team as a small community whose health can be judged in two ways.  
First, the people who will use or fund the product---\textbf{stakeholders}---should be happy with the result.  
Second, the members of the team themselves should feel motivated and satisfied with how they work together.  
When a team respects one another, admits that no one is infallible, and trusts each other to do the right work, daily interactions become smoother, decisions are made faster, and problems are solved cooperatively.  
A \textbf{selfmanaging team} takes these ideas further by giving the group the freedom to decide its own culture, task distribution, and quality standards, rather than following a topdown command.

\subsection{Formal Definition}
An \textbf{effective software team} is evaluated along two dimensions:  

\begin{itemize}
 \item \textbf{Stakeholder satisfaction}: Users or sponsors are pleased with the delivered outcome.  
 \item \textbf{Team satisfaction}: Team members feel motivated, valued, and content with their work environment.  
\end{itemize}

The \textbf{Three Pillars (HRT)} that support collaboration are:

\begin{enumerate}
 \item \textbf{Humility}: Team members recognize that they are not infallible, acknowledge room for improvement, and avoid equating personal identity with their code.  
 \item \textbf{Respect}: Feedback is directed at the work, not the person; members listen attentively and genuinely care about each other's contributions.  
 \item \textbf{Trust}: Individuals believe that others will act responsibly, deliver on commitments, and make fair decisions; this does not eliminate the need for review, but it establishes a baseline of confidence.  
\end{enumerate}

A \textbf{selfmanaging team} decides its own culture, allocates tasks, and sets quality practices without external command, leveraging autonomy to adapt quickly.

\textbf{Team leader} responsibilities include:

\begin{itemize}
 \item Acting as a catalyst rather than a boss, by setting clear goals and removing obstacles.
 \item Mediating conflicts and ensuring that nonperformers are addressed constructively.
 \item Supporting the HRT pillars and fostering both stakeholder and team satisfaction.
\end{itemize}

\textbf{Communication} occurs in two modes:

\begin{itemize}
 \item \textbf{Synchronous}: Short, regular standup meetings where team members speak in real time to share progress and raise blockers.
 \item \textbf{Asynchronous}: Issue trackers, comments, or emails that allow information to be exchanged without requiring simultaneous presence.
\end{itemize}

These definitions preempt common misunderstandings: agreement on every decision is not required; a single star programmer cannot replace a dysfunctional team; and trust does not mean unchecked work---review remains essential.

\subsection{Worked Example}
\textbf{Problem statement}\\
A fiveperson project consists of a product manager (PM), two developers (Dev A and Dev B), a tester, and a UX designer. During a sprint, a disagreement arises about how to implement a new filtering feature: Dev A proposes a quick hack, while Dev B argues for a more robust, reusable component. The team must resolve the conflict while maintaining stakeholder and team satisfaction.

\textbf{Step 1 -- Apply Humility}\\
Dev A acknowledges that the hack may cause future maintenance issues and says, ``I'm open to hearing a better approach; my solution isn't the only way.''  
Dev B likewise admits, ``My design will take longer now, but it will save time later. Let's weigh both options.''

\textbf{Step 2 -- Apply Respect}\\
During the standup (synchronous communication), each developer presents the pros and cons of their proposal without criticizing the other person. The tester adds, ``From a testing perspective, the reusable component reduces duplicate test cases.''  
The UX designer notes, ``A stable filter aligns with the UI flow we designed.''

\textbf{Step 3 -- Apply Trust}\\
The PM trusts both developers to act in the project's best interest and asks them to collaborate on a hybrid solution: implement the core of the reusable component now, and add a temporary wrapper for immediate release. The team agrees to a short timeline and commits to delivering the wrapper by the end of the day.

\textbf{Step 4 -- Use Asynchronous Communication}\\
Dev A documents the wrapper implementation in the issue tracker, tagging Dev B for review. Dev B adds comments explaining the reusable component's architecture and attaches design diagrams. The tester logs a test case in the tracker referring to the new wrapper. All updates are visible to the whole team without needing a meeting.

\textbf{Step 5 -- Review and Consolidate}\\
At the next standup, the team reviews the completed wrapper and discusses the plan to replace it with the full component in the next sprint. The PM reports to stakeholders that the feature is delivered on time, satisfying stakeholder expectations. Team members express satisfaction that their ideas were heard and integrated, boosting morale.

\textbf{Final answer}\\
The conflict was resolved through humility (admitting limits), respect (focusing on work, not people), and trust (relying on each other's commitment). Communication combined synchronous standup discussion with asynchronous issuetracker documentation, achieving both stakeholder satisfaction (feature delivered) and team satisfaction (collaborative process).

\textbf{What this shows us:}  
Applying the HRT pillars creates a respectful environment where disagreements become opportunities for better solutions, and the blend of synchronous and asynchronous communication keeps everyone aligned while preserving autonomy.

\subsection{Key Takeaway}
Effective software teams balance stakeholder satisfaction with team satisfaction, and they achieve this balance by practicing humility, respect, and trust within a selfmanaging, communicative environment.

\subsection{Questions}

\noindent \textbf{Recall 1.}  
According to the formal definition, what are the two dimensions used to evaluate whether a software team is effective?

\newpage

\noindent \textbf{Understand 1.}  
Explain in your own words how the pillar of \textbf{Respect} shapes the way team members give feedback to each other. Include why the feedback is directed at the work rather than the person.

\newpage

\noindent \textbf{Understand 2.}  
Provide a concrete example of how \textbf{Respect} can be demonstrated during a synchronous standup meeting when a disagreement arises over a technical solution.

\newpage

\noindent \textbf{Apply 1.}  
A developer on your team has missed three consecutive sprint deadlines. Using the \textbf{Trust} pillar, describe the steps you would take to address the situation while maintaining both stakeholder and team satisfaction.

\newpage

\noindent \textbf{Apply 2.}  
During a 10day sprint the team's stakeholdersatisfaction metric is 80 \% and the average teamsatisfaction survey score is 3.2 out of 5. The target for the next sprint is at least 85 \% stakeholder satisfaction and a teamsatisfaction score of 4.0.  
Calculate the additional percentagepoint increase required for each dimension to meet the targets.

\newpage

\noindent \textbf{Apply 3.}  
In a conflict similar to the worked example, Dev A wants a quick hack while Dev B proposes a reusable component. Outline, stepbystep, how the team leader can use the \textbf{Humility}, \textbf{Respect}, and \textbf{Trust} pillars to guide the team toward a collaborative solution.

\newpage

\noindent \textbf{Analyse 1.}  
Compare a \textbf{selfmanaging team} with a \textbf{traditional hierarchical team} with respect to (a) decisionmaking speed and (b) team morale. Highlight the advantages and disadvantages of each structure.

\newpage

\noindent \textbf{Analyse 2.}  
Identify a potential drawback of a selfmanaging team when its members have widely differing skill levels. Discuss how this drawback could affect the team's effectiveness and suggest one mitigation strategy.

\newpage

\noindent \textbf{Evaluate 1.}  
A team consistently delivers products that satisfy stakeholders (high stakeholdersatisfaction scores) but members report low morale and motivation (low teamsatisfaction scores). Assess whether such a team is likely to be sustainable over the long term. Provide at least two reasons supporting your judgment, referencing the concepts of the HRT pillars and selfmanaging teams.

\newpage
% ============================================================
% SECTION  Software Project Management & Effort Estimation
% ============================================================

\subsection*{Notation Used in This Section}
\begin{tabular}{@{}ll@{}}
\toprule
Symbol & Meaning \\ \midrule
    $Effort = a \times size^{b}$ & Personmonths effort as a function of software size (KLOC) with calibrated constants a and b \\
\bottomrule
\end{tabular}
\vspace{1em}

% !!!  REVIEW FLAG: Software Project Management & Effort Estimation  !!!
% [undefined_term] Add a brief definition in the theory section explaining that the planning phase involves estimating effort, defining scope, and creating project schedules.
% [undefined_term] Add a brief definition in the theory section explaining that the monitoring/control phase focuses on tracking progress, managing risks, and controlling changes to keep the project on schedule.
% !!! END REVIEW FLAGS !!!
\section{Software Project Management \& Effort Estimation}
\subsection{Before We Begin}
Regression analysis is a way of finding a mathematical relationship between two quantities by looking at a set of observed data.  For example, we might look at many past projects and see how the amount of code relates to the time the developers needed.  The Cone of Uncertainty describes how early guesses about a project's size or duration are very spread out, but as the design and requirements become clearer the spread becomes much smaller.  Both ideas help us understand why effort estimates change during a project's life.  With that background we can now examine how software projects estimate the work they will require.

\subsection{Motivation}
Estimating effort for a software project is like estimating how many bricks you need to build a house before you start buying them.  Project managers need a realistic picture of how much time and cost a project will consume so that they can plan resources, set schedules, and negotiate with customers.  Good estimates reduce the risk of overruns and help keep the team focused on delivering value.

\subsection{Intuition}
At its heart, effort estimation is about guessing how much work a team must perform to produce a certain amount of software.  Early in a project we have little information, so our guesses are broad; as we design, write requirements, and start coding, we learn more and can tighten those guesses.  Two common ways to make these guesses are: (1) start with an overall size and break it down (topdown), or (2) estimate each piece first and then add them together (bottomup).

\subsection{Formal Definition}
The basic topdown effort model used in many projectplanning tools is  

\[
\mbox{Effort}=a\times \mbox{size}^{\,b}
\]

where  

\begin{itemize}
\item \(\mbox{Effort}\) is the total work expressed in \textbf{personmonths} (one person working for one month).
\item \(\mbox{size}\) is the amount of software measured in \textbf{KLOC} (thousands of lines of source code).
\item \(a\) and \(b\) are constants that are calibrated from historical project data using regression analysis; they capture influences such as product reliability, complexity, and staff capability.
\end{itemize}

In practice, the COCOMO (Constructive Cost Model) family of models uses this formula but refines the estimate by multiplying by up to fifteen costdriver factors (e.g., reliability, analyst capability).  The constants \(a\) and \(b\) are \textit{not} universal; they differ for various types of software and development environments.  Likewise, bottomup estimation does not ignore risk---it simply builds the estimate from detailed componentlevel work items, which can then be adjusted for risk factors.  Neither approach is inherently more accurate; the choice depends on how much detailed information is already available.

\subsection{Worked Example}
\textbf{Problem statement:}  A project is expected to contain \(10\) KLOC of code.  Using the COCOMOtype formula with \(a = 2.5\) and \(b = 1.1\), calculate the estimated effort in personmonths.

\textbf{Step 1: Write the formula.}  
\[
\mbox{Effort}= 2.5\times (10)^{1.1}
\]

\textbf{Step 2: Compute the size term.}  
\(10^{1.1}=10^{1+0.1}=10\times10^{0.1}\).  
\(10^{0.1}\approx1.2589\).  
Thus \(10^{1.1}\approx10\times1.2589=12.589\).

\textbf{Step 3: Multiply by the constant \(a\).}  
\(\mbox{Effort}=2.5\times12.589\approx31.47\) personmonths.

\textbf{Final answer:}  The project's estimated effort is roughly 31.5 personmonths.

\textbf{What this shows us:}  Even with a simple sizeonly model, the formula translates a size measurement (KLOC) into a concrete work estimate, illustrating how topdown estimation can be performed quickly when only highlevel data are available.

\subsection{Key Takeaway}
Effort estimation links software size to required work via the formula \(\mbox{Effort}=a\times\mbox{size}^{b}\); understanding when to use topdown versus bottomup approaches lets managers produce realistic plans as project information matures.

\subsection*{Questions}
\textbf{Recall Question:}\\
List the two broad phases of software project management and give one sentence describing the main purpose of each phase.\\
\newpage
\textbf{Recall Question:}\\
Write the topdown effort estimation formula and define each of its symbols (\(\mbox{Effort}, a, b, \mbox{size}\)) including the units used.\\
\newpage
\textbf{Understanding Question:}\\
In your own words, explain why effort estimates tend to become more accurate as a project progresses from the feasibility stage to the coding stage. Refer to the ``Cone of Uncertainty'' concept in your answer.\\
\newpage
\textbf{Understanding Question:}\\
Describe how regression analysis is used to obtain the constants \(a\) and \(b\) in the effortsize formula. Why is historical project data essential for this step?\\
\newpage
\textbf{Application Question:}\\
Using the effortsize formula \(\mbox{Effort}=a \times \mbox{size}^{\,b}\) with \(a = 2.5\) and \(b = 1.1\), calculate the estimated effort for a project that is expected to contain \(8\) KLOC of code. Show all intermediate steps and round the final answer to one decimal place.\\
\newpage
\textbf{Application Question:}\\
A small maintenance project is estimated to contain \(5\) KLOC. The organization has calibrated the constants for such projects as \(a = 1.8\) and \(b = 1.2\).  
\begin{enumerate}
  \item Compute the estimated effort in personmonths.
  \item If the team can allocate 2 persons fulltime, estimate how many calendar months the effort will require (assume no overhead or holidays).
\end{enumerate}  
Provide calculations and the final rounded values.\\
\newpage
\textbf{Application Question:}\\
A project manager decides to use a bottomup approach. She breaks the system into three components with the following estimated efforts: Component A -- 10 personmonths, Component B -- 14 personmonths, Component C -- 9 personmonths.  
\begin{enumerate}
  \item What is the total effort estimate for the whole project using bottomup estimation?
  \item If a risk factor of 1.15 is applied to the total (to account for uncertainty), what is the adjusted effort estimate?
\end{enumerate}
Show the arithmetic used to reach each answer.\\
\newpage
\textbf{Analysis Question:}\\
Compare the advantages and disadvantages of topdown versus bottomup estimation for a project that currently has many unknown requirements. In your comparison, discuss at least two benefits and two limitations of each approach and indicate which situation would favour one method over the other.\\
\newpage
\textbf{Evaluation Question:}\\
Critique the statement: ``COCOMO provides exact effort predictions for any software project.''  
In your critique, consider the role of the constants \(a\) and \(b\), the influence of costdriver factors, the reliance on historical data, and the variability introduced by projectspecific characteristics. Provide a balanced argument that highlights both the strengths and the inherent limitations of the COCOMO model.\\
\newpage

\newpage
% ============================================================
% SECTION  Scheduling, Quality Planning, and Risk Management
% ============================================================

\subsection*{Notation Used in This Section}
\begin{tabular}{@{}ll@{}}
\toprule
Symbol & Meaning \\ \midrule
    $S = 2.5 \times E^{0.38}$ & Estimated overall schedule (months) as a function of effort (personmonths) \\
\bottomrule
\end{tabular}
\vspace{1em}

\section{Scheduling, Quality Planning, and Risk Management}

\subsection{Before We Begin}
When a software defect is discovered late in the lifecycle, fixing it often requires reworking many parts of the system, which makes the effort grow dramatically.  
A probabilityimpact risk matrix uses the words Low, Medium, and High on two axes to show how likely a risk is and how bad its consequences would be.  
Understanding these ideas helps us see why early defect removal matters and how to read a simple risk chart.  
With that picture in mind, we can now look at the ways projects plan their work and control quality.

\subsection{Motivation}
Finding and fixing a software defect early is like mopping up a small spill before it spreads across the entire floor.  
Project managers need tools that tell them how long a whole project will take and how to break that time into concrete tasks.  
They also need to know that catching defects sooner saves huge amounts of cost, and that some risks deserve immediate attention while others can be monitored.

\subsection{Intuition}
Think of a project schedule as a twolayer map.  
The overall schedule is the highlevel view that shows only the major milestones -- for example, design finished, implementation done, and testing complete.  
Below that, the detailed schedule lists every small activity, such as ``run unit tests for Module A,'' and assigns each activity to a team member and a specific amount of effort.  
Early defect removal works because a tiny mistake discovered while writing requirements is cheap to fix, while the same mistake later may require rewriting code, retesting, and even patching live systems.  
A risk matrix lets us spot the most dangerous spots: risks that are both very likely (High probability) and very damaging (High impact) are the ones we must act on first.

\subsection{Formal Definition}
The overall schedule \(S\) (measured in months) is estimated from the total effort \(E\) (measured in personmonths) by the empirical formula  

\[
S = 2.5 \times E^{0.38},
\]

where the symbol \(S = 2.5 \times E^{0.38}\) denotes this relationship.  

The detailed schedule consists of individual tasks; each task is defined by a duration (e.g., 1 day) and an effort amount (e.g., 7 hours).  The effort for a task is the actual persontime spent, not merely the calendar time.

Quality planning targets a low delivered defect density, expressed as the number of defects per thousand lines of code (defects/KLOC).  Defects may be introduced at any development stage, but they are removed through Quality Control activities such as reviews and testing.  The latency of a defect is the time between its introduction and its removal.  The cost to fix a defect grows with latency; fixing a defect in the operational phase can be about one hundred times more expensive than fixing it during the requirements phase.

In risk management, a risk is an uncertain event that could cause project failure.  Each risk receives a probability rating (Low, Medium, High) and an impact rating (Low, Medium, High).  Plotting these ratings on a twodimensional matrix creates a risk matrix.  Risks that fall in the HighProbability/HighImpact (HH) quadrant, or in the HighProbability/MediumImpact (HM) and MediumProbability/HighImpact (MH) quadrants, are prioritized for mitigation.  Mitigation may involve concrete steps such as rotating job assignments to guard against loss of trained personnel.

Note that the detailed schedule is not a mere copy of the overall schedule with finer numbers; it requires fresh planning of resources, dependencies, and effort estimates. Likewise, not all defects cost the same to fix---the later a defect is discovered, the higher its fixing cost.  Finally, risks that are highprobability but lowimpact still merit attention because they can accumulate or combine with other issues.

\subsection{Worked Example}
\textbf{Problem statement.}  
A project is estimated to require \(E = 20\) personmonths of effort.  Compute the overall schedule \(S\) using the formula, then divide the schedule into three milestones (design, implementation, testing) and show one detailed task (unit testing) with its duration and effort.

\medskip
\textbf{Step 1: Compute the overall schedule.}  
\[
S = 2.5 \times 20^{0.38}.
\]  
First evaluate the exponent: \(20^{0.38} \approx 3.121\).  
Multiply by 2.5: \(S \approx 2.5 \times 3.121 = 7.80\) months.  

\textbf{Step 2: Allocate milestones.}  
Assume a practical split of the overall time:  
\begin{itemize}
  \item Design: 30 % of \(S\) \(= 0.30 \times 7.80 \approx 2.34\) months,
  \item Implementation: 50 % of \(S\) \(= 0.50 \times 7.80 \approx 3.90\) months,
  \item Testing: 20 % of \(S\) \(= 0.20 \times 7.80 \approx 1.56\) months.
\end{itemize}

\textbf{Step 3: Define a detailed task.}  
During the testing milestone, one concrete activity is ``Unit test Module X.''  
\begin{itemize}
  \item Duration: 1 calendar day,
  \item Effort: 7 hours (which is \(7/8 = 0.875\) persondays, assuming an 8hour workday).
\end{itemize}

\textbf{Final answer.}  
The project's overall schedule is about 7.8 months.  The milestones are roughly 2.3 months (design), 3.9 months (implementation), and 1.6 months (testing).  A sample detailed task for the testing phase is ``Unit test Module X,'' lasting 1 day and requiring 7 hours of effort.

\textbf{What this shows us.}  
The formula gives a quick, toplevel estimate of total time, while the detailed schedule breaks that time into concrete activities that teams can assign and track.  Both layers are needed for effective project control.

\subsection{Key Takeaway}
The overall schedule provides a highlevel time estimate using \(S = 2.5 \times E^{0.38}\), and the detailed schedule refines that estimate into specific tasks that enable early defect removal and focused risk mitigation.

\subsection*{Questions}
\newpage
\textbf{1. Recall.}  
Write the empirical relationship that estimates the overall project schedule \(S\) (in months) from the total effort \(E\) (in personmonths).

\newpage
\textbf{2. Understand.}  
In your own words, explain why fixing a defect during the requirements phase costs far less than fixing the same defect after the software has been deployed and is operating in production.

\newpage
\textbf{3. Understand.}  
Describe the purpose of a probabilityimpact risk matrix. How does placing risks on the HighProbability/HighImpact (HH) and the adjacent HM/MH quadrants help a project manager decide where to focus mitigation effort?

\newpage
\textbf{4. Apply.}  
A project is estimated to require \(E = 12\) personmonths of effort.  
Using the formula \(S = 2.5 \times E^{0.38}\), calculate the approximate overall schedule \(S\) (in months). Show your intermediate steps.

\newpage
\textbf{5. Apply.}  
Using the overall schedule \(S\) you obtained in the previous question, allocate the time to three milestones --- Design (30 % of \(S\)), Implementation (50 % of \(S\)), and Testing (20 % of \(S\)).  
State the duration of each milestone in months (round to two decimal places).

\newpage
\textbf{6. Analyse.}  
Consider the following three risks:

\begin{enumerate}
  \item \textbf{Risk A}: Probability = High, Impact = Medium,
  \item \textbf{Risk B}: Probability = Medium, Impact = High,
  \item \textbf{Risk C}: Probability = Low, Impact = High.
\end{enumerate}

Place each risk in the appropriate cell of a probabilityimpact matrix (Low, Medium, High on both axes) and construct a priority list from highest to lowest urgency for mitigation.  Explain the reasoning behind the ordering.

\newpage
\textbf{7. Evaluate.}  
Assess the statement: ``A detailed schedule adds little value if the overall schedule is already known.''  
Provide a justified argument drawing on the concepts of overall vs. detailed scheduling, early defect removal, and risk mitigation discussed in this section.

\newpage
% ============================================================
% SECTION  Software Process Models
% ============================================================

% !!!  REVIEW FLAG: Software Process Models  !!!
% [undefined_term] Introduce and define "defect density" in the formal definition or earlier discussion.
% !!! END REVIEW FLAGS !!!
\section{Software Process Models}
\subsection{Motivation}
Choosing a process model is like choosing a route to travel from your home to a new city---some routes are straight highways (waterfall), others are stopandgo buses (iterative).  Different routes help teams reach the destination efficiently under different conditions.  Understanding the strengths and weaknesses of each model lets you match a project's uncertainty, risk, and schedule to the most suitable path.  This section explains the major models and how they relate to effort and defect distribution.

\subsection{Intuition}
A software process is simply the ordered set of activities a team follows to turn ideas into a working program.  Some teams move linearly from requirements to deployment, while others loop back, build quick sketches, or fix a time limit for each cycle.  The key is to pick the pattern that reduces risk, handles changing requirements, and fits the amount of time and resources available.

\subsection{Formal Definition}
A \textbf{software process} is a collection of steps that produce software.  The \textbf{distribution of effort} across the lifecycle typically allocates roughly 10--20 \% of total effort to Requirements, 10--20 \% to Design, 20--30 \% to Coding, and 30--50 \% to Testing.  Although \textbf{Coding} exhibits the highest \textbf{defect density} (over 50 \% of defects originate there), it is not the most expensive phase because testing consumes the largest share of effort and resources.

The four major process models are:
\begin{itemize}
  \item \textbf{Waterfall}: A strictly linear sequence of phases -- Requirements $\rightarrow$ HighLevel Design $\rightarrow$ Detailed Design $\rightarrow$ Coding $\rightarrow$ Testing $\rightarrow$ Deployment.  Each phase finishes before the next begins, producing extensive documentation.  It works well when the problem is well understood and requirements are unlikely to change.  A common mistake is to think the waterfall model can easily accommodate frequent requirement changes; in reality, late changes are costly.
  \item \textbf{Prototyping}: The team quickly builds a \emph{prototype}---a simplified, often ``quick and dirty'' version of the system---to explore and clarify requirements.  Feedback from users refines the requirements before full development starts.  It is important not to assume the prototype becomes the final product; typically the prototype is discarded or heavily reworked.
  \item \textbf{Iterative Development}: The overall project is divided into a series of miniwaterfalls, each delivering a functional increment.  Lessons from one iteration feed into the next, avoiding unnecessary features and allowing gradual improvement.  While this reduces feature bloat, an unplanned architecture can suffer if the overall structure is not considered early.
  \item \textbf{Timeboxing}: A fixed length of time (the ``box'') is assigned to an iteration, and the team selects the subset of features that fit inside that box.  If the scope is too large, some features are postponed.  The schedule is therefore predictable, and stages such as design, coding, and testing can be pipelined to run in parallel.  Timeboxing does not mean sacrificing quality; quality must be built into the work that fits within the box.
\end{itemize}

\subsection{Worked Example}
\textbf{Problem statement:} A small team is commissioned to create a mobileapp that displays local restaurant menus.  The client is unsure about the exact layout and interaction style for the user interface (UI).  Determine which process model should be used to gather requirements, and then decide whether iterative development or timeboxing is more appropriate for the later development phases.

\textbf{Step 1 -- Identify the main uncertainty.}  
The uncertainty lies in the UI design; the functional backend (fetching menus) is well understood.

\textbf{Step 2 -- Choose a model to reduce this uncertainty.}  
Because the goal is to explore and clarify UI requirements, the \textbf{prototyping} model is ideal.  A quick mockup of the app screens can be shown to users, whose feedback will shape the final UI specifications.

\textbf{Step 3 -- Evaluate the next phase options.}  
After the UI requirements are solidified, the project still faces two constraints: (a) the client wants a usable product within six weeks, and (b) the team wants to deliver value early.

\textbf{Step 4 -- Compare iterative development and timeboxing.}  
\begin{itemize}
  \item \textbf{Iterative development} would break the remaining work into several increments (e.g., basic menu display, search feature, user reviews).  Each increment would be a miniwaterfall, allowing the team to incorporate feedback from earlier increments.
  \item \textbf{Timeboxing} would fix the sixweek schedule and select the most critical features that fit within that window (e.g., core menu display and simple search).  Less critical features would be postponed to a later release.
\end{itemize}

\textbf{Step 5 -- Decide the best fit.}  
Because the overall schedule is tight (six weeks) and the team already has clear UI requirements, \textbf{timeboxing} is the better choice for the followup phases.  It guarantees delivery of a functional core app within the deadline while allowing the team to pipeline design, coding, and testing.

\textbf{Final answer:} Use prototyping to clarify UI requirements, then adopt timeboxing for the remaining development to meet the fixed schedule.

\textbf{What this shows us:} Selecting a process model starts by matching the model's strength to the project's biggest uncertainty; once that uncertainty is resolved, the scheduledriven nature of timeboxing can ensure timely delivery, whereas iterative development is preferable when the timeline is flexible and continuous feedback is desired.

\subsection{Key Takeaway}
Choose the process model whose main characteristic (linear, exploratory, incremental, or schedulefixed) best addresses the project's dominant uncertainty and constraints.

\subsection{Questions}
\newpage
\textbf{Recall 1.} List the four software process models that were described in this section.

\newpage
\textbf{Recall 2.} Which software process model is best suited when requirements have low likelihood of change?

\newpage
\textbf{Understand 1.} In your own words, explain why the coding phase, even though it generates more than 50 of defects, is not the most expensive phase of a software project.

\newpage
\textbf{Understand 2.} Describe how the distribution of effort across the lifecycle (requirements, design, coding, testing) relates to the overall cost of a software project.

\newpage
\textbf{Apply 1.} A small team is starting a project whose requirements are welldefined and the likelihood of change is low.  Which software process model should they choose for the development work, and what characteristic of that model makes it suitable for this situation?

\newpage
\textbf{Apply 2.} Suppose the total estimated effort for a project is 1,200 personhours.  
Using the effort distribution given in the section (Requirements 10--20 , Design 10--20 , Coding 20--30 , Testing 30--50 ), calculate a possible allocation of personhours to each phase that respects the indicated ranges and sums to the total effort.

\newpage
\textbf{Analyse 1.} Compare the way risk is handled in the prototyping model versus the timeboxing model.  Which types of risk are each model better at mitigating, and why?

\newpage
\textbf{Analyse 2.} In the worked example, the team first used prototyping to resolve UI uncertainty and then chose timeboxing for the remaining work.  Explain how this sequence manages both requirements uncertainty and schedule risk, and discuss an alternative sequencing (e.g., prototyping followed by iterative development) and its possible impact on risk.

\newpage
\textbf{Evaluate 1.} Critically evaluate the claim: ``Iterative development always results in higher quality software than waterfall.''  Consider the strengths and weaknesses of each model presented in the theory, and discuss conditions under which the claim might be true or false.

\newpage
% ============================================================
% SECTION  Open Source Software (OSS) Processes
% ============================================================

\section{Open Source Software (OSS) Processes}

\subsection{Before We Begin}
A versioncontrol system stores the full history of a software project in a place called a repository.  
When a developer makes a change, that change is recorded as a commit, which captures the new state of the files.  
A patch is a compact description of a single change that can be applied to the repository to produce the same result as a commit.  
Projects often work with one central repository so that everyone can see and build on the same baseline code.  
Understanding these ideas will help you follow how opensource software moves from a user's request to a released feature.  

\section{Motivation}
\emph{An OSS project is like a community garden where anyone can bring seeds, but experienced gardeners decide which seeds become part of the main plot.}  
Opensource software (OSS) exists because many people can contribute ideas, fixes, and new features without a single buyer dictating the work.  
This collaborative model solves the problem of limited resources in a single company by pooling the effort of volunteers and organizations around the world.  

\section{Intuition}
Think of an OSS project as a living ecosystem that starts with a small prototype created by a handful of programmers (the Conception Stage).  
When that prototype proves useful, a broader group of developers, testers, and users joins, forming a Community Stage where anyone may suggest improvements.  
Roles are loosely defined: a core developer writes most of the code, a peripheral contributor helps occasionally, and a user reports problems or ideas.  
Work flows from the user's request, through volunteerwritten patches, to review by core developers, and finally into an official release managed by a release manager.  

\section{Formal Definition}
An opensource software (OSS) project proceeds through two lifecycle stages:

\begin{itemize}
  \item \textbf{Conception Stage}: A few programmers create an initial working prototype.
  \item \textbf{Community Stage}: The prototype is opened to a globally distributed community that communicates asynchronously.
\end{itemize}

The community contains several roles:

\begin{itemize}
  \item \textbf{Project Leader}: Guides overall direction but does not alone decide which code is accepted.
  \item \textbf{Core Developers} (about 15 people in typical projects): Write roughly eighty percent of new features and twothirds of bug fixes; they also review and approve patches.
  \item \textbf{Peripheral Contributors}: Volunteer developers who pick issues from a public log, create patches, and submit them for review.
  \item \textbf{Users}: Report problems or request new functionality; their reports become \emph{Modification Requests (MRs)} in an issuetracking system such as Bugzilla.
\end{itemize}

The workflow for handling an MR is:

\begin{enumerate}
  \item A user files an MR describing a desired change or bug fix.
  \item A peripheral contributor selects the MR, writes a \emph{patch} (a text description of the code change), and submits it to the issue tracker.
  \item A core developer reviews the patch, possibly requesting revisions, and then \emph{commits} the change to the \emph{central repository} (the single authoritative copy of the source code, often managed with CVS or a similar system).
  \item The \emph{release manager} periodically freezes the set of accepted features, stabilizes the build, and creates a new released version.
\end{enumerate}

Two important ideas counter common misunderstandings:

\begin{itemize}
  \item \textit{Not only core developers submit code.} Peripheral contributors regularly provide patches, and their contributions can be merged after review.
  \item \textit{Quality control exists despite volunteer effort.} Core developers enforce review standards, and the release manager ensures that only vetted changes reach a release.
  \item \textit{Users influence development.} Their MRs drive what issues the community works on, shaping the project's direction.
\end{itemize}

\textbf{Linus's Law} states that ``given enough eyeballs, all bugs are shallow,'' meaning that the large number of participants increases the chance that any defect will be noticed and fixed.  This principle explains why many OSS projects achieve high reliability even though contributions are voluntary.  

\subsection{Worked Example}
\textbf{Problem statement}  
In the fictitious OSS project \emph{LibreCalc}, a user reports a lowpriority bug: ``The ``Save'' button does not respond when the file name contains a space.''  The bug is entered as a Modification Request (MR101) in the project's Bugzilla system.  Follow the typical OSS workflow from this MR to its inclusion in the next release.

\begin{enumerate}
  \item \textbf{User files the MR} -- The user creates MR101, describing the observed behavior, the environment, and a minimal reproducible example.  The MR is now visible to the entire community.
  \item \textbf{Peripheral contributor selects the MR} -- Alice, a peripheral contributor, scans the open MR list, sees MR101, and decides to fix it. She checks out the latest code from the central repository.
  \item \textbf{Alice creates a patch} -- Alice edits the source code to correctly handle file names with spaces, tests the change locally, and generates a patch file named \texttt{fixsavespace.patch}.  The patch shows the exact lines added and removed.
  \item \textbf{Alice submits the patch} -- Alice attaches \texttt{fixsavespace.patch} to MR101 in Bugzilla and adds a short comment explaining her solution.
  \item \textbf{Core developer reviews the patch} -- Bob, a core developer, receives a notification about the new attachment. He reviews the code, runs the full test suite, and finds that the patch fixes the bug without side effects.  He writes a brief review comment approving the change.
  \item \textbf{Core developer commits the patch} -- Bob applies the patch to his local copy of the repository and creates a commit with the message ``Fix Save button failure when file name contains spaces (MR101).''  The commit is pushed to the central CVS repository, making the fix part of the official code base.
  \item \textbf{Release manager freezes the next build} -- Carol, the release manager, announces a feature freeze for the upcoming version 1.2. She verifies that MR101 is now part of the repository, includes it in the build, runs integration tests, and tags the release.
\end{enumerate}

The bug reported by the user is now fixed, merged into the central repository, and shipped in LibreCalc version 1.2.  The example demonstrates the complete OSS workflow: a usergenerated MR, volunteerwritten patch, coredeveloper review and commit, and finally inclusion by the release manager.  It also illustrates the distributed, collaborative nature of responsibility and quality control in opensource projects.

\subsection{Key Takeaway}
Opensource projects move from a user's request to a released feature through a transparent, communitydriven process where anyone can propose changes, core developers ensure quality, and release managers deliver stable versions.

\subsection*{Questions}
\begin{enumerate}
  \item \textbf{Recall.} List and briefly describe the two lifecycle stages that an opensource software (OSS) project typically passes through, according to the theory section.
  \newpage
  \item \textbf{Understand.} In your own words, trace the path of a Modification Request (MR) from the moment a user files it until the change is merged into the central repository. Mention the key roles and actions involved at each step.
  \newpage
  \item \textbf{Understand.} Identify which participant (or group of participants) is responsible for each of the following activities in the OSS workflow:  
  \begin{enumerate}
    \item Creating a patch that implements an MR.  
    \item Reviewing the patch and deciding whether it can be committed.  
    \item Freezing the code base for an upcoming release.
  \end{enumerate}
  \newpage
  \item \textbf{Apply.} You have discovered a lowpriority bug in an OSS project and wish to fix it as a peripheral contributor. Outline, in sequential order, the concrete steps you would take---from selecting the MR to having your change incorporated into the next official release. Be sure to include any artefacts (e.g., patch file, commit message) you would need to produce.
  \newpage
  \item \textbf{Apply.} Assume the bug you are fixing concerns a ``Save'' button that fails when the file name contains a space (similar to the LibreCalc example).  
  \begin{enumerate}
    \item Write an appropriate filename for the patch you would generate.  
    \item Draft a concise commit message that follows the style shown in the worked example.
  \end{enumerate}
  \newpage
  \item \textbf{Analyse.} Compare and contrast the motivations of an individual (extrinsic vs. intrinsic) contributor with those of a corporate contributor to an OSS project. How might these differing motivations influence the type of work each chooses to do?
  \newpage
  \item \textbf{Evaluate.} Critically assess the claim: ``Linus's Law guarantees no bugs will ever remain in an OSS project.'' Discuss whether the law alone is sufficient to achieve a bugfree product, citing specific mechanisms or safeguards described in the theory.
\end{enumerate}
\newpage

\newpage
% ============================================================
% SECTION  Agile Development
% ============================================================

% !!!  REVIEW FLAG: Agile Development  !!!
% [undefined_term] Define user stories before using them, e.g., as concise, usercentered descriptions of desired functionality.
% [undefined_term] Introduce the Planning Game as part of XP's release planning process before referencing it in the formal definitions.
% !!! END REVIEW FLAGS !!!
\section{Agile Development}

\subsection{Before We Begin}
Refactoring means changing the inner structure of a piece of code while keeping what the code does visible to users exactly the same.  
It is done after a test that proves the current behavior is correct, so the developer can be sure the changes have not broken anything.  
Think of it as reorganizing the furniture in a room without moving any of the walls.  
The purpose is to make the code easier to understand, extend, and maintain.  
With that idea in mind, we can now look at how Agile and Extreme Programming use short cycles and continuous improvement.

\paragraph{Motivation}
\textit{Analogy first:} Agile development is like preparing a series of small tasting plates rather than cooking a huge banquet all at once.  
This approach exists because trying to deliver a massive, fully documented product in one go often leads to wasted effort when requirements change.  
By delivering short, working pieces of software and constantly talking with the customer, teams can adapt quickly and keep value flowing.

\paragraph{Intuition}
Imagine a kitchen where chefs prepare a bitesize dish, taste it, get feedback from the diners, and then immediately improve the recipe before moving on to the next bite.  
In software, Agile teams build a tiny, usable feature, show it to the customer, learn what works and what does not, and then tweak the next piece accordingly.  
The cycle repeats many times, producing a complete meal (the final product) that matches what the diners actually want.

\paragraph{Formal Definition}
\begin{itemize}
  \item \textbf{Agile Manifesto Values} -- the four guiding statements are:
  \begin{enumerate}
    \item Working software over comprehensive documentation;
    \item Customer collaboration over contract negotiation;
    \item Responding to change over following a plan;
    \item Individuals and interactions over processes and tools.
  \end{enumerate}
  \item \textbf{Extreme Programming (XP) ProjectManagement Practices}
  \begin{enumerate}
    \item Small, crossfunctional teams work together in an open office; the team keeps the same members (continuity) and invites an onsite customer representative.
    \item The \emph{Planning Game} creates a release plan from user stories (highlevel requirements) ordered by risk and value, then an iteration plan that breaks each story into small tasks.
    \item Daily standup meetings let each person share what was done, what will be done, and any obstacles.
    \item The team respects a 40hour work week to avoid burnout.
  \end{enumerate}
  \item \textbf{XP Technical Practices}
  \begin{enumerate}
    \item \emph{User stories} capture what the customer needs in simple language.
    \item \emph{TestDriven Development (TDD)} follows a threestep cycle: write a failing test (red), write just enough code to make the test pass (green), then improve the code's structure without changing its behavior (refactor).
    \item \emph{Continuous Integration (CI)} requires developers to check in code at least once a day; an automated build verifies that the whole system still compiles and passes all tests.
    \item \emph{Pair Programming} has two developers share one workstation, constantly reviewing each other's work; this improves code quality without necessarily doubling the time needed.
    \item \emph{Simple incremental design} and \emph{refactoring} keep the design as simple as possible while the system grows.
    \item \emph{Shared code} means the whole team owns every part of the code base, so anyone can improve any piece.
  \end{enumerate}
\end{itemize}
These definitions avoid the misconception that pair programming always slows delivery, that TDD means writing all tests before any code runs, or that Agile eliminates all documentation.

\paragraph{Worked Example}
\textbf{Goal:} Create a function \verb|factorial(n)| that returns the product of all positive integers up to \verb|n|, using TDD, and illustrate pair programming and continuous integration.

\begin{enumerate}
  \item \textbf{Step 1 -- Write a failing test (Red).}
        \begin{verbatim}
        // Test case 1
        assertEquals(1, factorial(0));   // 0! = 1
        // Test case 2
        assertEquals(1, factorial(1));   // 1! = 1
        // Test case 3
        assertEquals(2, factorial(2));   // 2! = 2
        // Test case 4
        assertEquals(6, factorial(3));   // 3! = 6
        \end{verbatim}
        The test suite is run; it fails because \verb|factorial| does not exist yet.

  \item \textbf{Step 2 -- Write the minimal code to make the test pass (Green).}
        \begin{verbatim}
        int factorial(int n) {
            return 1;   // temporary implementation
        }
        \end{verbatim}
        Running the tests now yields failure for cases where \verb|n| > 1, because the implementation always returns 1.

        The team (two developers sitting together in pair programming) discusses the simplest change that will satisfy the next failing test. They replace the body with a small conditional:
        \begin{verbatim}
        int factorial(int n) {
            if (n == 0) return 1;
            if (n == 1) return 1;
            if (n == 2) return 2;
            if (n == 3) return 6;
            return 0;   // placeholder for larger n
        }
        \end{verbatim}
        All four tests now pass. The pair commits the code to the shared repository; the CI server automatically builds the project and runs the full test suite, confirming green status.

  \item \textbf{Step 3 -- Refactor the implementation (Refactor).}
        The pair notices the code is repetitive and hard to extend. They agree on a clean, recursive definition:
        \begin{verbatim}
        int factorial(int n) {
            if (n <= 1) return 1;
            return n  factorial(n - 1);
        }
        \end{verbatim}
        No new tests are needed because the existing suite already guarantees correctness for the values covered. The refactored code is checked in; the CI system rebuilds and reruns all tests, which still pass.

  \item \textbf{Step 4 -- Extend the test suite for a larger value (Red).}
        \begin{verbatim}
        // Additional test
        assertEquals(24, factorial(4));   // 4! = 24
        \end{verbatim}
        Running the tests now fails (the recursive method works, but the test reveals a bug in handling larger numbers due to integer overflow if using a small type; the pair decides to change the return type to a larger numeric type).

  \item \textbf{Step 5 -- Adjust code to satisfy the new test (Green).}
        \begin{verbatim}
        long factorial(long n) {
            if (n <= 1) return 1L;
            return n  factorial(n - 1);
        }
        \end{verbatim}
        After changing the method signature and updating all calls, the CI run reports all tests green again.

  \item \textbf{Final answer.}
        The function \verb|factorial| now correctly computes factorials for the tested values, the code is concise, and the whole team has continuously integrated their changes, receiving immediate feedback from the automated build.
\end{enumerate}
\textbf{What this shows us:} The TDD cycle (red--green--refactor) lets a team create just enough code to satisfy a test, then improve it without breaking anything, while pair programming and continuous integration provide rapid, shared quality checks throughout the development process.

\subsection{Key Takeaway}
Agile and Extreme Programming combine short, feedbackrich iterations with disciplined practices such as TDD, pair programming, and continuous integration to deliver highquality software that evolves safely and rapidly.

\subsection{Questions}
\textbf{Question 1 (Recall).}\\
List the four values that appear in the Agile Manifesto.\\
\newpage
\textbf{Question 2 (Understand).}\\
Describe each of the three steps of the TestDriven Development (TDD) cycle (Red, Green, Refactor). For each step, explain why it is required in order to keep the code correct while it evolves.\\
\newpage
\textbf{Question 3 (Understand).}\\
In your own words, explain how the practice of pair programming contributes to code quality and how it interacts with Continuous Integration (CI) during a development iteration.\\
\newpage
\textbf{Question 4 (Apply).}\\
A user story reads: ``As a user, I want to sort a list of integers in ascending order so that I can view the data ordered.''*  
Outline, step by step, how an XP team would use pair programming and Continuous Integration to deliver this story within a single iteration. Your answer should include:
\begin{enumerate}
  \item How the story is broken down into tasks.
  \item How the pair designs, implements, and tests the sorting feature.
  \item How the code is checked in and how CI validates the change.
\end{enumerate}
\\
\newpage
\textbf{Question 5 (Apply).}\\
During the same iteration, the team decides to adopt the TDD approach for the sorting functionality.  
Write a short example of a failing test (Red) for sorting the list \{3,1,2\}, then sketch the minimal implementation (Green) that would make this test pass, and finally describe one refactoring (Refactor) you would perform to improve the code without changing its behaviour.\\
\newpage
\textbf{Question 6 (Analyze).}\\
Compare daily standup meetings with a weekly status report. Discuss how daily standups help an XP team identify and mitigate risks earlier, and why a weekly report might miss or delay detection of the same risks.\\
\newpage
\textbf{Question 7 (Evaluate).}\\
Some critics claim that XP's strong emphasis on continuous refactoring creates unnecessary code churn, slowing delivery. Evaluate this argument. In your answer, provide at least two reasons why regular refactoring is beneficial for an agile project, and consider any potential downsides that the critics highlight. Your evaluation should be grounded in the definitions and practices described in the section.\\
\newpage

\end{document}
